\documentclass[10pt,a4paper]{article}
\usepackage{fontspec}
\setmainfont{Times New Roman}
\setsansfont{Arial}
\setmonofont{Consolas}
\usepackage{amsmath,amssymb}
\usepackage{xcolor}
\usepackage{tabularx}
\usepackage{array}
\usepackage{multicol}

% Page geometry using native TeX dimensions
\setlength{\textwidth}{182mm}
\setlength{\textheight}{254mm}
\setlength{\oddsidemargin}{-11mm}
\setlength{\evensidemargin}{-11mm}
\setlength{\topmargin}{-18mm}
\setlength{\headheight}{0mm}
\setlength{\headsep}{0mm}
\setlength{\footskip}{10mm}
\setlength{\parindent}{0pt}
\setlength{\parskip}{2.5pt}

% Color definitions
\definecolor{primary}{RGB}{15, 45, 95}
\definecolor{defblue}{RGB}{238, 245, 255}
\definecolor{defframe}{RGB}{25, 75, 155}
\definecolor{thmgreen}{RGB}{240, 250, 240}
\definecolor{thmframe}{RGB}{20, 115, 40}
\definecolor{propamber}{RGB}{255, 251, 235}
\definecolor{propframe}{RGB}{175, 115, 15}
\definecolor{proofbg}{RGB}{247, 247, 247}
\definecolor{proofframe}{RGB}{115, 115, 115}

% Custom styled block commands
\newcommand{\DefBox}[2]{%
\noindent\fcolorbox{defframe}{defblue}{%
\parbox{\dimexpr\textwidth-2\fboxsep-2\fboxrule\relax}{%
\textbf{\color{defframe}#1}\par\vspace{1.5pt}#2%
}%
}\vspace{2.5pt}%
}

\newcommand{\ThmBox}[2]{%
\noindent\fcolorbox{thmframe}{thmgreen}{%
\parbox{\dimexpr\textwidth-2\fboxsep-2\fboxrule\relax}{%
\textbf{\color{thmframe}#1}\par\vspace{1.5pt}#2%
}%
}\vspace{2.5pt}%
}

\newcommand{\PropBox}[2]{%
\noindent\fcolorbox{propframe}{propamber}{%
\parbox{\dimexpr\textwidth-2\fboxsep-2\fboxrule\relax}{%
\textbf{\color{propframe}#1}\par\vspace{1.5pt}#2%
}%
}\vspace{2.5pt}%
}

\newcommand{\ProofBox}[1]{%
\noindent\fcolorbox{proofframe}{proofbg}{%
\parbox{\dimexpr\textwidth-2\fboxsep-2\fboxrule\relax}{%
\textit{\color{proofframe}Краткое доказательство:}\par\vspace{1.5pt}#1 \hfill $\blacksquare$%
}%
}\vspace{3.5pt}%
}

\begin{document}

% Title Header
\begin{center}
{\huge\textbf{\color{primary}ДИСКРЕТНАЯ МАТЕМАТИКА}}\\[4pt]
{\Large\textbf{Полный академический свод определений, теорем, свойств и доказательств}}\\[3pt]
{\large\textbf{Материалы к коллоквиуму (Лекции 1--11 + Логика первого порядка)}}\\[4pt]
{\normalsize Университет ИТМО $\cdot$ ФИТиП $\cdot$ Осенний семестр 2026 $\cdot$ Автор курса: Константин Чухарев}
\end{center}
\vspace{-4pt}
\hrule height 1.5pt
\vspace{4pt}

\tableofcontents
\vspace{6pt}
\hrule
\vspace{6pt}

\section{Логика высказываний (Лекция 1)}

\DefBox{Определение 1.1 (Высказывание)}{%
\textbf{Высказывание} --- это повествовательное утверждение, о котором можно однозначно сказать, истинно оно ($1$, $\mathrm{True}$) или ложно ($0$, $\mathrm{False}$). Вопросы, команды, пожелания и парадоксальные самоприменимые суждения высказываниями не являются.
}

\PropBox{Свойство 1.1 (Вакуумная истинность / Vacuous Truth)}{%
Импликация с ложной посылкой тождественно истинна при любом заключении:
\[ 0 \to Q \equiv 1 \quad (\forall Q). \]
Условие «если $P$, то $Q$» обещает выполнение $Q$ лишь при наступлении $P$. Если $P = 0$, обещание не нарушено. В теории множеств именно поэтому пустое множество является подмножеством любого множества ($\varnothing \subseteq A$), так как формула $\forall x \, (x \in \varnothing \to x \in A)$ вакуумно истинна.
}

\DefBox{Определение 1.2 (Логические связки и приоритеты)}{%
\begin{enumerate}
\item \textbf{Отрицание} ($\neg A$, $\overline{A}$): меняет значение на противоположное ($\neg 0 = 1, \neg 1 = 0$).
\item \textbf{Конъюнкция} ($A \land B$): истинна тогда и только тогда, когда $A = 1$ и $B = 1$.
\item \textbf{Дизъюнкция} ($A \lor B$): истинна, если хотя бы один из операндов равен $1$.
\item \textbf{Импликация} ($A \to B$): ложна исключительно при $1 \to 0$; во всех остальных случаях истинна ($A \to B \equiv \neg A \lor B$).
\item \textbf{Эквивалентность} ($A \leftrightarrow B$): истинна тогда и только тогда, когда $A = B$.
\end{enumerate}
\textbf{Приоритет связок (по убыванию старшинства):} $\neg \succ \land \succ \lor \succ \to \succ \leftrightarrow$. Скобки имеют абсолютный приоритет. Импликация правоассоциативна: $A \to B \to C \equiv A \to (B \to C)$.
}

\DefBox{Определение 1.3 (Интерпретация и модель)}{%
Пусть $V = \{p_1, \dots, p_n\}$ --- множество пропозициональных переменных. \textbf{Интерпретацией} называется функция $\nu: V \to \{0, 1\}$. Интерпретация гомоморфно продолжается на все формулы по таблицам истинности связок. Интерпретация $\nu$, при которой $\nu(F) = 1$, называется \textbf{моделью формулы} $F$ ($\nu \models F$). Для $n$ переменных существует ровно $2^n$ интерпретаций.
}

\DefBox{Определение 1.4 (Классификация формул по истинности)}{%
\begin{itemize}
\item \textbf{Тавтология (общезначимая):} истинна при всех интерпретациях ($\models F$). Пример: $A \lor \neg A$.
\item \textbf{Противоречие (тождественно ложная / невыполнимая):} ложна при всех интерпретациях. Пример: $A \land \neg A$.
\item \textbf{Выполнимая формула:} истинна хотя бы при одной интерпретации.
\item \textbf{Контингентная (нейтральная):} имеет и модели, и опровержения.
\end{itemize}
Критерий: $\models F \iff \neg F \text{ невыполнима (противоречие)}$.
}

\DefBox{Определение 1.5 (Логическая эквивалентность)}{%
Формулы $A$ и $B$ называются \textbf{логически эквивалентными} ($A \equiv B$), если их истинностные значения совпадают при всех интерпретациях: $\forall \nu: \nu(A) = \nu(B)$. Критерий: $A \equiv B \iff \models (A \leftrightarrow B)$.
}

\ThmBox{Теорема 1.1 (Законы де Моргана)}{%
\[ \neg (A \land B) \equiv \neg A \lor \neg B, \qquad \neg (A \lor B) \equiv \neg A \land \neg B. \]
}
\ProofBox{%
Проверяется таблицей истинности: левая часть $\neg(A \land B)$ ложна только при $A=1, B=1$. Правая часть $\neg A \lor \neg B$ также обращается в 0 только при $\neg A = 0$ и $\neg B = 0$, то есть при $A=1, B=1$. Значения совпадают во всех 4 строках.
}

\ThmBox{Теорема 1.2 (Отрицание и выражение импликации)}{%
\[ A \to B \equiv \neg A \lor B, \qquad \neg(A \to B) \equiv A \land \neg B. \]
}
\ProofBox{%
$A \to B$ ложна только при $A=1, B=0$. Формула $\neg A \lor B$ при $A=1$ даёт $0 \lor B = B$, что равно 0 только при $B=0$; при $A=0$ даёт $1 \lor B = 1$. Таблицы истинности совпадают. Отрицание: $\neg(A \to B) \equiv \neg(\neg A \lor B) \equiv \neg\neg A \land \neg B \equiv A \land \neg B$ по закону де Моргана и двойного отрицания.
}

\ThmBox{Теорема 1.3 (Закон контрапозиции)}{%
\[ A \to B \equiv \neg B \to \neg A. \]
}
\ProofBox{%
По выражению импликации: $\neg B \to \neg A \equiv \neg(\neg B) \lor \neg A \equiv B \lor \neg A \equiv \neg A \lor B \equiv A \to B$. Контрапозиция лежит в основе доказательства от противного.
}

\PropBox{Свойство 1.2 (Законы алгебры логики высказываний)}{%
\begin{itemize}
\item \textbf{Коммутативность:} $A \land B \equiv B \land A, \quad A \lor B \equiv B \lor A$.
\item \textbf{Ассоциативность:} $(A \land B) \land C \equiv A \land (B \land C), \quad (A \lor B) \lor C \equiv A \lor (B \lor C)$.
\item \textbf{Дистрибутивность:} $A \land (B \lor C) \equiv (A \land B) \lor (A \land C), \quad A \lor (B \land C) \equiv (A \lor B) \land (A \lor C)$.
\item \textbf{Идемпотентность:} $A \land A \equiv A, \quad A \lor A \equiv A$.
\item \textbf{Поглощение:} $A \land (A \lor B) \equiv A, \quad A \lor (A \land B) \equiv A$.
\item \textbf{Константы 0 и 1:} $A \land 1 \equiv A, \, A \land 0 \equiv 0, \, A \lor 0 \equiv A, \, A \lor 1 \equiv 1$.
\item \textbf{Дополнение:} $A \land \neg A \equiv 0, \quad A \lor \neg A \equiv 1, \quad \neg\neg A \equiv A$.
\item \textbf{Экспортация:} $(A \land B) \to C \equiv A \to (B \to C)$.
\end{itemize}
}

\DefBox{Определение 1.6 (Логическое следствие)}{%
Формула $B$ называется \textbf{логическим следствием} формул $A_1, \dots, A_n$ ($A_1, \dots, A_n \models B$), если при любой интерпретации $\nu$, где $\nu(A_1) = \dots = \nu(A_n) = 1$, заключение также истинно: $\nu(B) = 1$.
}

\ThmBox{Теорема 1.4 (Семантическая теорема о дедукции для высказываний)}{%
\[ A_1, \dots, A_n \models B \iff \models (A_1 \land \dots \land A_n) \to B. \]
}
\ProofBox{%
1. ($\implies$) Пусть $A_1, \dots, A_n \models B$. Рассмотрим произвольную $\nu$. Если все $\nu(A_i) = 1$, то по условию $\nu(B) = 1 \implies 1 \to 1 = 1$. Если хотя бы одна $\nu(A_k) = 0$, то посылка равна 0, и по вакуумной истинности $0 \to \nu(B) = 1$. Формула истинна всегда ($\models$).
2. ($\impliedby$) Пусть импликация --- тавтология. На любой строке, где $\nu(A_1) = \dots = \nu(A_n) = 1$, посылка равна 1. Чтобы импликация была истинна, заключение $\nu(B)$ не может быть 0, значит $\nu(B) = 1$. Следование доказано.
}

\PropBox{Свойство 1.3 (Фундаментальные правила вывода)}{%
\begin{enumerate}
\item \textbf{Modus Ponens (правило отделения):} $P \to Q, \, P \models Q$. Тавтология: $((P \to Q) \land P) \to Q \equiv 1$.
\item \textbf{Modus Tollens (правило отрицания):} $P \to Q, \, \neg Q \models \neg P$. Тавтология: $((P \to Q) \land \neg Q) \to \neg P \equiv 1$.
\end{enumerate}
}

\section{Предикаты и кванторы (Лекция 2)}

\DefBox{Определение 2.1 (Предикат и предметная область)}{%
\textbf{Предикат} $P(x_1, \dots, x_n)$ --- высказывательная форма, истинность которой зависит от подстановки конкретных значений переменных из \textbf{предметной области (домена / domain of discourse)}. Сам по себе предикат не истинен и не ложен.
}

\DefBox{Определение 2.2 (Кванторы $\forall$ и $\exists$)}{%
\begin{itemize}
\item \textbf{Квантор всеобщности ($\forall$):} $\forall x \, P(x)$ истинно $\iff$ для каждого объекта $a$ предметной области $P(a) = 1$ (бесконечная конъюнкция $P(a_1) \land P(a_2) \land \dots$). Ложно при наличии хотя бы одного контрпримера.
\item \textbf{Квантор существования ($\exists$):} $\exists x \, P(x)$ истинно $\iff$ найдётся хотя бы один свидетель $a$, для которого $P(a) = 1$ (бесконечная дизъюнкция $P(a_1) \lor P(a_2) \lor \dots$).
\end{itemize}
}

\DefBox{Определение 2.3 (Свободные и связанные переменные, предложения)}{%
Переменная под действием квантора называется \textbf{связанной}. Не связанная квантором переменная называется \textbf{свободной}. Формула без свободных переменных называется \textbf{замкнутой формулой (предложением)} и выражает конкретное высказывание ($0$ или $1$). Формула со свободными переменными называется открытой и задаёт предикат.
}

\ThmBox{Теорема 2.1 (Законы де Моргана для кванторов)}{%
\[ \neg(\forall x \, P(x)) \equiv \exists x \, \neg P(x), \qquad \neg(\exists x \, P(x)) \equiv \forall x \, \neg P(x). \]
}
\ProofBox{%
1. $\neg(\forall x \, P(x)) = 1 \iff$ суждение «для всех $x$ верно $P(x)$» ложно $\iff$ найдётся объект $x_0$, для которого $P(x_0) = 0 \iff \neg P(x_0) = 1 \iff \exists x \, \neg P(x) = 1$.
2. $\neg(\exists x \, P(x)) = 1 \iff$ в домене нет ни одного $x$ со свойством $P(x) \iff$ каждый $x$ удовлетворяет $\neg P(x) \iff \forall x \, \neg P(x) = 1$.
}

\PropBox{Свойство 2.1 (Порядок кванторов)}{%
\begin{itemize}
\item \textbf{Разноименные кванторы не коммутируют:} $\exists y \, \forall x \, P(x, y) \implies \forall x \, \exists y \, P(x, y)$, но обратное в общем случае неверно. ($\forall x \exists y$ разрешает выбор $y$, зависящий от $x$; $\exists y \forall x$ требует общего $y$ сразу для всех $x$).
\item \textbf{Одноименные кванторы перестановочны:} $\forall x \forall y \, P(x, y) \equiv \forall y \forall x \, P(x, y)$ и $\exists x \exists y \, P(x, y) \equiv \exists y \exists x \, P(x, y)$.
\end{itemize}
}

\DefBox{Определение 2.4 (Ограниченные кванторы и квантор единственности)}{%
\begin{itemize}
\item \textbf{Ограниченная всеобщность через импликацию:} $\forall x \in A \, P(x) \stackrel{\mathrm{def}}{\iff} \forall x \, (x \in A \to P(x))$. При $A = \varnothing$ посылка ложна $\implies \forall x \in \varnothing \, P(x) \equiv 1$ (вакуумная истинность).
\item \textbf{Ограниченное существование через конъюнкцию:} $\exists x \in A \, P(x) \stackrel{\mathrm{def}}{\iff} \exists x \, (x \in A \land P(x))$. При $A = \varnothing$ формула тождественно ложна: $\exists x \in \varnothing \, P(x) \equiv 0$.
\item \textbf{Квантор единственности ($\exists!$):} $\exists! x \, P(x) \stackrel{\mathrm{def}}{\iff} \exists x \, \Big( P(x) \land \forall y \, (P(y) \to y = x) \Big)$.
\end{itemize}
}

\section{Методы математических доказательств (Лекция 3)}

\DefBox{Определение 3.1 (Математическое доказательство)}{%
\textbf{Доказательство} --- проверяемая конечная цепочка логических рассуждений, выводящая утверждение из аксиом, определений и ранее доказанных теорем по строгим правилам логического вывода.
}

\PropBox{Свойство 3.1 (Классификация методов доказательства)}{%
1. \textbf{Прямое доказательство $P \to Q$:} предполагаем $P = 1$ и дедукцией получаем $Q = 1$.\\
2. \textbf{Доказательство через контрапозицию:} доказываем эквивалентную импликацию $\neg Q \to \neg P$.\\
3. \textbf{Доказательство от противного (Reductio ad absurdum):} предполагаем $\neg P$ и выводим противоречие вида $B \land \neg B$ или ложь $\bot$. Следовательно, $\neg P$ ложно, то есть истинно $P$.\\
4. \textbf{Разбор случаев:} если область разбита на покрытие $C_1 \lor \dots \lor C_k$, доказывается $\bigwedge_{i=1}^k (C_i \to Q)$.\\
5. \textbf{Контрпример:} для опровержения $\forall x \, P(x)$ предъявляется ровно один объект $x_0$ со свойством $\neg P(x_0)$.
}

\ThmBox{Теорема 3.1 (Чётность квадратов)}{%
Целое число $n$ чётно тогда и только тогда, когда $n^2$ чётно.
}
\ProofBox{%
($\implies$, прямое) Пусть $n = 2k$ ($k \in \mathbb{Z}$). Тогда $n^2 = (2k)^2 = 4k^2 = 2(2k^2) = 2m$, где $m = 2k^2 \in \mathbb{Z} \implies n^2$ чётно.
($\impliedby$, контрапозиция) Докажем «если $n$ нечётно, то $n^2$ нечётно». Пусть $n = 2k + 1$. Тогда $n^2 = (2k+1)^2 = 4k^2 + 4k + 1 = 2(2k^2 + 2k) + 1 = 2m + 1$, нечётно. По закону контрапозиции из чётности $n^2$ следует чётность $n$.
}

\ThmBox{Теорема 3.2 (Иррациональность $\sqrt{2}$)}{%
Число $\sqrt{2}$ иррационально: $\sqrt{2} \notin \mathbb{Q}$.
}
\ProofBox{%
Предположим противное: $\sqrt{2} \in \mathbb{Q}$, то есть $\sqrt{2} = \frac{a}{b}$, где $a, b \in \mathbb{Z}, b \ne 0, \gcd(a, b) = 1$ (несократимая дробь).
Возводя в квадрат: $2 = \frac{a^2}{b^2} \implies a^2 = 2b^2 \implies a^2$ чётно $\implies a$ чётно (по теореме 3.1).
Пусть $a = 2k$. Тогда $(2k)^2 = 2b^2 \implies 4k^2 = 2b^2 \implies b^2 = 2k^2 \implies b^2$ чётно $\implies b$ чётно.
Следовательно, и $a$, и $b$ делятся на 2, что противоречит $\gcd(a, b) = 1$. Противоречие $\implies \sqrt{2} \notin \mathbb{Q}$.
}

\ThmBox{Теорема 3.3 (Мультипликативность модуля)}{%
Для любых действительных $x, y$: $|x| \cdot |y| = |xy|$.
}
\ProofBox{%
Разбор случаев: 1) Если $x=0$ или $y=0$, то $0 = 0$. 2) Если $x, y > 0$, то $|x||y| = xy = |xy|$. 3) Если $x, y < 0$, то $|x||y| = (-x)(-y) = xy = |xy|$ (т.к. $xy > 0$). 4) Если разных знаков ($x > 0, y < 0$), то $|x||y| = x(-y) = -xy = |xy|$ (т.к. $xy < 0 \implies |xy| = -xy$). Все случаи разобраны.
}

\ThmBox{Теорема 3.4 (Иррациональная степень иррационального числа)}{%
Существуют такие иррациональные числа $a$ и $b$, что $a^b \in \mathbb{Q}$.
}
\ProofBox{%
(Неконструктивное доказательство) Рассмотрим число $x = \sqrt{2}^{\sqrt{2}}$. Число $\sqrt{2}$ иррационально.
1) Если $x \in \mathbb{Q}$, то полагаем $a = \sqrt{2}, b = \sqrt{2}$, и теорема доказана.
2) Если $x \notin \mathbb{Q}$, то положим $a = x = \sqrt{2}^{\sqrt{2}}$ (иррационально) и $b = \sqrt{2}$ (иррационально). Тогда $a^b = (\sqrt{2}^{\sqrt{2}})^{\sqrt{2}} = \sqrt{2}^{\sqrt{2} \cdot \sqrt{2}} = \sqrt{2}^2 = 2 \in \mathbb{Q}$. В обоих случаях такие числа существуют.
}

\ThmBox{Теорема 3.5 (Принцип математической индукции)}{%
Пусть $P(n)$ --- предикат на $\mathbb{N}$. Если выполнена база $P(1) = 1$, и для любого $k \ge 1$ верен индукционный шаг $P(k) \implies P(k+1)$, то $P(n) = 1$ для всех $n \in \mathbb{N}$.
Формула суммы: $1 + 2 + \dots + n = \frac{n(n+1)}{2}$. Сумма кубов: $1^3 + \dots + n^3 = (1 + \dots + n)^2$.
}
\ProofBox{%
База для суммы ($n=1$): $1 = \frac{1 \cdot 2}{2} = 1$. Шаг ($k \to k+1$): $1 + \dots + k + (k+1) = \frac{k(k+1)}{2} + (k+1) = (k+1)(\frac{k}{2} + 1) = \frac{(k+1)(k+2)}{2}$. Формула доказана для всех $n$.
}

\ThmBox{Теорема 3.6 (Сильная индукция и основная теорема арифметики)}{%
Если $P(n_0) = 1$ и $\forall k \ge n_0 : (\forall i \le k : P(i)) \implies P(k+1)$, то $P(n) = 1$ для всех $n \ge n_0$.
Всякое целое $n \ge 2$ раскладывается в произведение простых чисел.
}
\ProofBox{%
База: $n=2$ простое. Шаг: пусть верно для всех $i \le k$. Для $k+1$: если $k+1$ простое --- доказано. Если составное, то $k+1 = a \cdot b$, где $2 \le a, b \le k$. По предположению сильной индукции $a$ и $b$ раскладываются в простые множители, значит и их произведение $k+1 = a \cdot b$ раскладывается.
}

\ThmBox{Теорема 3.7 (Принцип вполне упорядоченности $\mathbb{N}$)}{%
Любое непустое подмножество $S \subseteq \mathbb{N}$ содержит наименьший элемент: $\exists m \in S \; \forall x \in S : m \le x$. Принцип вполне упорядоченности логически эквивалентен принципу математической индукции.
}

\section{Логика множеств и наивная теория множеств (Лекция 4)}

\DefBox{Определение 4.1 (Универсум и способы задания множеств)}{%
\textbf{Универсум $\mathcal{U}$} --- объемлющее фиксированное множество всех допустимых объектов.
1. \textbf{Экстенсиональный (перечисление):} $A = \{x_1, \dots, x_n\}$.
2. \textbf{Интенсиональный (характеристический предикат):} $A = \{x \in \mathcal{U} \mid P(x)\}$. Индикатор принадлежности: $[x \in A] \in \{0, 1\}$.
}

\PropBox{Свойство 4.1 (Парадокс Рассела)}{%
Принцип неограниченного свёртывания Кантора приводит к противоречию: рассмотрим множество всех множеств, не содержащих себя: $R = \{x \mid x \notin x\}$. Тогда $R \in R \iff R \notin R$ (тупик).
Следствие: универсального множества всех множеств не существует. В аксиоматике ZFC выделение подмножеств разрешено только из уже построенного множества.
}

\DefBox{Определение 4.2 (Включение, равенство и булеан)}{%
\begin{itemize}
\item \textbf{Подмножество:} $A \subseteq B \iff \forall x \, (x \in A \to x \in B)$.
\item \textbf{Строгое подмножество:} $A \subsetneq B \iff A \subseteq B \land A \ne B$.
\item \textbf{Равенство множеств (метод двух включений):} $A = B \iff (A \subseteq B) \land (B \subseteq A)$.
\item \textbf{Булеан (Power Set):} множество всех подмножеств $\mathcal{P}(A) = 2^A = \{X \mid X \subseteq A\}$.
\end{itemize}
}

\ThmBox{Теорема 4.1 (Мощность булеана конечного множества)}{%
Если $|A| = n$, то $|\mathcal{P}(A)| = 2^n$. Число подмножеств размера $k$ равно $\binom{n}{k}$.
}
\ProofBox{%
Каждое подмножество $X \subseteq A$ взаимно-однозначно кодируется двоичным вектором длины $n$: $(\chi_X(a_1), \dots, \chi_X(a_n)) \in \{0, 1\}^n$. Число таких векторов равно $2 \times 2 \times \dots \times 2 = 2^n$. Число способов выбрать ровно $k$ единиц среди $n$ позиций равно числу сочетаний $\binom{n}{k}$.
}

\ThmBox{Теорема 4.2 (Законы де Моргана и поглощения для множеств)}{%
\[ \overline{A \cup B} = \overline{A} \cap \overline{B}, \quad \overline{A \cap B} = \overline{A} \cup \overline{B}, \qquad A \cup (A \cap B) = A, \quad A \cap (A \cup B) = A. \]
}
\ProofBox{%
Для поглощения $A \cup (A \cap B) = A$: 1) $A \subseteq A \cup (A \cap B)$ тривиально. 2) Если $x \in A \cup (A \cap B)$, то либо $x \in A$, либо $x \in A \cap B \implies x \in A$. В обоих случаях $x \in A \implies A \cup (A \cap B) \subseteq A$. По двум включениям множества равны.
}

\PropBox{Свойство 4.2 (Симметрическая разность и индикаторы)}{%
Симметрическая разность $A \triangle B = (A \setminus B) \cup (B \setminus A)$. Характеристическая функция:
\[ \chi_{A \triangle B}(x) = \chi_A(x) \oplus \chi_B(x) = (\chi_A(x) + \chi_B(x)) \bmod 2. \]
Элемент $x \in A_1 \triangle \dots \triangle A_k \iff x$ принадлежит нечётному числу множеств $A_i$.
}

\DefBox{Определение 4.3 (Пара Куратовского и декартово произведение)}{%
\textbf{Упорядоченная пара:} $(a, b) \stackrel{\mathrm{def}}{=} \{\{a\}, \{a, b\}\}$.\\
\textbf{Декартово произведение:} $A \times B = \{(a, b) \mid a \in A \land b \in B\}$. Для конечных множеств: $|A \times B| = |A| \cdot |B|$.
}

\ThmBox{Теорема 4.3 (Фундаментальное свойство упорядоченной пары)}{%
\[ (a, b) = (c, d) \iff a = c \;\land\; b = d. \]
}
\ProofBox{%
Пусть $\{\{a\}, \{a, b\}\} = \{\{c\}, \{c, d\}\}$.
Пересечение элементов: $\bigcap (a, b) = \{a\} \cap \{a, b\} = \{a\}$. Аналогично $\bigcap (c, d) = \{c\}$. Значит, $\{a\} = \{c\} \implies a = c$.
Объединение элементов: $\bigcup (a, b) = \{a, b\}$.
Если $a = b$, то $(a, a) = \{\{a\}\}$, откуда $\{c, d\} = \{a\} \implies c = d = a = b$.
Если $a \ne b$, то из $\{a, b\} = \{c, d\}$ при $a = c$ следует $b = d$.
}

\section{Бинарные отношения и эквивалентность (Лекция 5)}

\DefBox{Определение 5.1 (Бинарное отношение и представления)}{%
\textbf{Бинарное отношение} между $A$ и $B$ --- любое подмножество $R \subseteq A \times B$. Если $A=B$, то $R \subseteq A^2$ --- отношение на $A$. Запись: $(a, b) \in R \iff aRb$.
Представления: пары, булева матрица $M_R \in \{0, 1\}^{n \times n}$ ($M_{ij} = 1 \iff a_i R a_j$), ориентированный граф (дуга $a \to b \iff aRb$).
}

\PropBox{Свойство 5.1 (Базовые свойства отношений на $A$)}{%
\begin{itemize}
\item \textbf{Рефлексивность:} $\forall a: aRa$ ($M_{ii} = 1$, петли у всех вершин).
\item \textbf{Иррефлексивность:} $\forall a: \neg(aRa)$ ($M_{ii} = 0$, полное отсутствие петель).
\item \textbf{Симметричность:} $\forall a, b: aRb \implies bRa$ ($M = M^T$).
\item \textbf{Антисимметричность:} $\forall a, b: (aRb \land bRa) \implies a = b$ (нет встречных дуг между разными вершинами).
\item \textbf{Асимметричность:} $\forall a, b: aRb \implies \neg(bRa)$ (нет петель и нет встречных дуг).
\item \textbf{Транзитивность:} $\forall a, b, c: (aRb \land bRc) \implies aRc$ ($M_R \odot M_R \le M_R$).
\end{itemize}
}

\ThmBox{Теорема 5.1 (Критерий асимметричности)}{%
Отношение асимметрично тогда и только тогда, когда оно иррефлексивно и антисимметрично.
}
\ProofBox{%
1. ($\implies$) Если $aRa$, то по асимметричности $\neg(aRa)$ (противоречие) $\implies$ иррефлексивно. Если $aRb \land bRa$, то при $a \ne b$ асимметричность нарушена $\implies a = b$ (антисимметрично).
2. ($\impliedby$) Пусть $aRb$. Если бы $bRa$, то по антисимметричности $a = b \implies aRa$, что противоречит иррефлексивности. Значит, $\neg(bRa)$.
}

\DefBox{Определение 5.2 (Операции и композиция отношений)}{%
Обратное отношение: $R^{-1} = \{(b, a) \mid (a, b) \in R\}$. Матрица: $M_{R^{-1}} = M_R^T$. $R$ симметрично $\iff R = R^{-1}$.\\
Композиция: $S \circ R = \{(a, c) \mid \exists b : (a, b) \in R \land (b, c) \in S\}$. Матрица: $M_{S \circ R} = M_R \odot M_S$.
}

\ThmBox{Теорема 5.2 (Свойства композиции)}{%
1. Ассоциативность: $(T \circ S) \circ R = T \circ (S \circ R)$.\\
2. Обращение: $(S \circ R)^{-1} = R^{-1} \circ S^{-1}$.\\
3. Транзитивность через композицию: $R$ транзитивно $\iff R \circ R \subseteq R$.
}
\ProofBox{%
Для п. 3: $(a, c) \in R \circ R \iff \exists b: aRb \land bRc$. Если $R$ транзитивно, отсюда $aRc \implies (a, c) \in R$, значит $R \circ R \subseteq R$. Обратно: если $aRb$ и $bRc$, то $(a, c) \in R \circ R \subseteq R \implies aRc$.
}

\DefBox{Определение 5.3 (Замыкания отношений)}{%
Замыкание отношения $R$ относительно свойства $\mathcal{P}$ --- наименьшее отношение $R' \supseteq R$, обладающее свойством $\mathcal{P}$.
Рефлексивное: $R^= = R \cup \mathrm{id}_A$. Симметричное: $R^s = R \cup R^{-1}$. Транзитивное: $R^+ = \bigcup_{k=1}^n R^k$. Рефлексивно-транзитивное: $R^* = R^+ \cup \mathrm{id}_A$.
}

\ThmBox{Теорема 5.3 (Алгоритм Уоршелла, $\mathcal{O}(n^3)$)}{%
Пусть $W^{(0)} = M_R$. При последовательном добавлении промежуточных вершин $k = 1, \dots, n$:
\[ W^{(k)}[i][j] = W^{(k-1)}[i][j] \;\lor\; \Big( W^{(k-1)}[i][k] \;\land\; W^{(k-1)}[k][j] \Big). \]
Матрица $W^{(n)}$ в точности равна матрице транзитивного замыкания $R^+$.
}

\DefBox{Определение 5.4 (Отношение эквивалентности и фактор-множество)}{%
Отношение $\sim$ называется \textbf{эквивалентностью}, если оно рефлексивно, симметрично и транзитивно.
\textbf{Класс эквивалентности:} $[a]_\sim = \{x \in A \mid x \sim a\}$. \textbf{Фактор-множество:} $A/\sim \;= \{[a]_\sim \mid a \in A\}$.
}

\ThmBox{Теорема 5.4 (Фундаментальные свойства классов эквивалентности)}{%
1. $a \in [a]$.\\
2. $[a] = [b] \iff a \sim b$.\\
3. Либо $[a] = [b]$, либо $[a] \cap [b] = \varnothing$ (классы не имеют частичных пересечений).
}
\ProofBox{%
1. По рефлексивности $a \sim a \implies a \in [a]$.
2. ($\implies$) Если $[a]=[b]$, то $a \in [b] \implies a \sim b$. ($\impliedby$) Пусть $a \sim b$. Если $x \in [a]$, то $x \sim a \land a \sim b \implies x \sim b \implies x \in [b]$. Значит $[a] \subseteq [b]$. По симметрии $[b] \subseteq [a] \implies [a] = [b]$.
3. Пусть $[a] \cap [b] \ne \varnothing$, т.е. существует $c \in [a] \cap [b]$. Тогда $c \sim a$ и $c \sim b \implies a \sim c \land c \sim b \implies a \sim b \implies [a] = [b]$.
}

\ThmBox{Теорема 5.5 (Связь эквивалентностей и разбиений)}{%
Всякая эквивалентность $\sim$ задаёт разбиение $A$ на непустые непересекающиеся блоки $\{[a]\}$. Обратно, любое разбиение $\mathcal{P} = \{A_i\}$ задаёт отношение эквивалентности: $a \sim b \iff \exists i : a, b \in A_i$.
}

\ThmBox{Теорема 5.6 (Корректность операций на фактор-множестве / Построение $\mathbb{Q}$)}{%
На $P = \mathbb{Z} \times (\mathbb{Z} \setminus \{0\})$ отношение $(a, b) \sim (c, d) \iff ad = bc$ является эквивалентностью. Поле рациональных чисел $\mathbb{Q} = P/\sim$. Операция сложения $[(a, b)] + [(c, d)] := [(ad + bc, bd)]$ корректна (не зависит от выбора представителей).
}
\ProofBox{%
Пусть $(a', b') \sim (a, b)$ ($a'b = b'a$) и $(c', d') \sim (c, d)$ ($c'd = d'c$). Проверим $(a'd' + b'c')bd \stackrel{?}{=} b'd'(ad + bc)$.
LHS $= (a'b)(d'd) + (c'd)(b'b) = (b'a)(d'd) + (d'c)(b'b) = b'd'(ad + bc) =$ RHS. Равенство верно.
}

\section{Отношения порядка, ЧУМ и решётки (Лекция 6)}

\DefBox{Определение 6.1 (Частичный и линейный порядок, ЧУМ)}{%
\begin{itemize}
\item \textbf{Нестрогий частичный порядок ($\preceq$):} рефлексивен, антисимметричен, транзитивен. Пара $(A, \preceq)$ --- ЧУМ (Poset).
\item \textbf{Строгий порядок ($\prec$):} иррефлексивен и транзитивен (автоматически асимметричен).
\item \textbf{Линейный (тотальный) порядок:} любые два элемента сравнимы ($\forall a, b: a \preceq b \lor b \preceq a$).
\end{itemize}
}

\ThmBox{Теорема 6.1 (Асимметричность строгого порядка и взаимный переход)}{%
1. Всякий строгий порядок асимметричен: $a \prec b \implies \neg(b \prec a)$.\\
2. Нестрогий порядок порождает строгий: $a \prec b \iff a \preceq b \land a \ne b$. Строгий порождает нестрогий: $a \preceq b \iff a \prec b \lor a = b$.
}
\ProofBox{%
Если бы $a \prec b$ и $b \prec a$, то по транзитивности $a \prec a$, что противоречит иррефлексивности.
}

\DefBox{Определение 6.2 (Цепи, антицепи, ширина и высота)}{%
\textbf{Цепь} $C \subseteq A$ --- подмножество попарно сравнимых элементов. \textbf{Антицепь} $S \subseteq A$ --- подмножество попарно несравнимых элементов ($x \ne y \implies x \not\preceq y \land y \not\preceq x$).
\textbf{Высота} $\mathrm{height}(A) = \max |C|$. \textbf{Ширина} $\mathrm{width}(A) = \max |S|$ (мера максимального параллелизма независимых задач).
}

\ThmBox{Теорема 6.2 (Лемма о спуске в конечном ЧУМ)}{%
Всякое конечное непустое ЧУМ содержит хотя бы один \textbf{минимальный} и хотя бы один \textbf{максимальный} элемент.
}
\ProofBox{%
Если бы минимального элемента не было, для любого $a_0$ существовал бы $a_1 \prec a_0$, затем $a_2 \prec a_1$ и т.д. В силу конечности элементы повторятся: $a_i = a_j$ ($i < j$). По транзитивности $a_i \prec a_i$, противоречие с иррефлексивностью.
}

\ThmBox{Теорема 6.3 (Существование топологической сортировки / Алгоритм Кана)}{%
Любое конечное ЧУМ (DAG) допускает согласованный линейный порядок: $a \preceq b \implies a \le b$.
Алгоритм Кана находит сортировку за время $\mathcal{O}(|V| + |E|)$.
}
\ProofBox{%
По лемме о спуске выбираем минимальный элемент $m_1$, ставим его первым и удаляем. В оставшемся конечном множестве снова выбираем минимальный элемент $m_2$ и т.д. За $|V|$ шагов все элементы упорядочены в цепь без нарушения зависимостей.
}

\DefBox{Определение 6.3 (Отношение покрытия, диаграммы Хассе, экстремальные элементы)}{%
\textbf{Покрытие:} $a \prec: b \iff a \prec b$ и $\nexists c : a \prec c \prec b$. \textbf{Диаграмма Хассе:} транзитивное сокращение орграфа (только рёбра покрытия $\prec:$, снизу вверх, без петель).\\
\textbf{Минимальный:} $\nexists x : x \prec m$ (локальное свойство, может быть много). \textbf{Наименьший:} $\forall x : m_0 \preceq x$ (глобальное свойство).
}

\ThmBox{Теорема 6.4 (Единственность наименьшего/наибольшего элемента)}{%
Наименьший (наибольший) элемент ЧУМ, если существует, то строго единственен. Всякий наименьший элемент минимален.
}
\ProofBox{%
Пусть $m_1, m_2$ --- наименьшие. Тогда $m_1 \preceq m_2$ (т.к. $m_1$ наименьший) и $m_2 \preceq m_1$ (т.к. $m_2$ наименьший). По антисимметричности $m_1 = m_2$.
}

\DefBox{Определение 6.4 (Супремум, инфимум, решётки)}{%
Точная верхняя грань: $\sup S = \bigvee S$ (наименьшая из верхних граней). Точная нижняя грань: $\inf S = \bigwedge S$ (наибольшая из нижних граней).\\
ЧУМ $(L, \preceq)$ называется \textbf{решёткой (Lattice)}, если для любых $a, b \in L$ существуют $a \vee b = \sup(a, b)$ и $a \wedge b = \inf(a, b)$.
}

\ThmBox{Теорема 6.5 (Дистрибутивные и модулярные решётки / Критерии Дедекинда и Биркгофа)}{%
1. Решётка дистрибутивна: $a \wedge (b \vee c) = (a \wedge b) \vee (a \wedge c)$.\\
2. Решётка модулярна: $a \preceq c \implies a \vee (b \wedge c) = (a \vee b) \wedge c$.\\
3. \textbf{Критерий Дедекинда:} решётка модулярна $\iff$ не содержит пентагон $N_5$ в качестве подрешётки.\\
4. \textbf{Критерий Биркгофа:} решётка дистрибутивна $\iff$ не содержит ни ромб $M_3$, ни пентагон $N_5$. Всякая дистрибутивная решётка модулярна.
}

\ThmBox{Теорема 6.6 (Кнастера --- Тарского о неподвижной точке)}{%
Всякая монотонная функция $f: L \to L$ на полной решётке $(L, \preceq)$ имеет наименьшую неподвижную точку $\mathrm{lfp}(f)$ и наибольшую неподвижную точку $\mathrm{gfp}(f)$:
\[ \mathrm{lfp}(f) = \bigwedge \{x \in L \mid f(x) \preceq x\}, \qquad \mathrm{gfp}(f) = \bigvee \{x \in L \mid x \preceq f(x)\}. \]
}
\ProofBox{%
Пусть $S = \{x \in L \mid f(x) \preceq x\}$ и $m = \bigwedge S$. Для любого $s \in S$: $m \preceq s \implies f(m) \preceq f(s) \preceq s$. Значит, $f(m)$ --- нижняя грань для $S \implies f(m) \preceq \bigwedge S = m$. Применяя монотонность: $f(f(m)) \preceq f(m)$, откуда $f(m) \in S$. Тогда $m \preceq f(m)$ (т.к. $m$ нижняя грань $S$). Из $f(m) \preceq m$ и $m \preceq f(m)$ следует $f(m) = m$.
}

\ThmBox{Теорема 6.7 (Теорема Дилуорса и теорема Мирского)}{%
В любом конечном ЧУМ:
\[ \min |\text{разбиение на цепи}| = \max |\text{антицепь}| \quad (\text{Дилуорс, 1950}), \]
\[ \min |\text{разбиение на антицепи}| = \max |\text{цепь}| \quad (\text{Мирский, 1971}). \]
}

\section{Функции, отображения и асимптотический анализ (Лекция 7)}

\DefBox{Определение 7.1 (Функция как отношение, частичные функции)}{%
Функция $f: A \to B$ --- это бинарное отношение $f \subseteq A \times B$, удовлетворяющее:
1) \textbf{Тотальности:} $\forall a \in A \; \exists b \in B : (a, b) \in f$;
2) \textbf{Функциональности:} $(a, b_1) \in f \land (a, b_2) \in f \implies b_1 = b_2$.
Число функций $|B^A| = |B|^{|A|} = m^n$. Если тотальность не требуется, функция называется \textbf{частичной} ($f: A \rightharpoonup B$).
Критерий графика: любая вертикальная прямая пересекает график тотальной функции ровно один раз.
}

\DefBox{Определение 7.2 (Инъекция, сюръекция, биекция)}{%
\begin{itemize}
\item \textbf{Инъекция:} $\forall a_1, a_2 : a_1 \ne a_2 \implies f(a_1) \ne f(a_2)$ ($\iff f(a_1) = f(a_2) \implies a_1 = a_2$).
\item \textbf{Сюръекция:} $\forall b \in B \; \exists a \in A : f(a) = b \iff f(A) = B$.
\item \textbf{Биекция:} инъекция и сюръекция одновременно ($\forall b \in B \; \exists! a \in A : f(a) = b$).
\end{itemize}
}

\ThmBox{Теорема 7.1 (Критерии обратимости функций)}{%
1. $f: A \to B$ инъективна $\iff$ существует левая обратная $g: B \to A$ ($g \circ f = \mathrm{id}_A$).\\
2. $f: A \to B$ сюръективна $\iff$ существует правая обратная $g: B \to A$ ($f \circ g = \mathrm{id}_B$, аксиома выбора).\\
3. $f: A \to B$ биективна $\iff$ существует единственная двусторонняя обратная $f^{-1}: B \to A$.
}
\ProofBox{%
Для п. 3 (единственность): пусть $g_1, g_2$ --- две обратные. Тогда $g_1 = g_1 \circ \mathrm{id}_B = g_1 \circ (f \circ g_2) = (g_1 \circ f) \circ g_2 = \mathrm{id}_A \circ g_2 = g_2$.
}

\ThmBox{Теорема 7.2 (Свойства композиции)}{%
1. Композиция сохраняет тип: $f, g$ инъективны $\implies g \circ f$ инъективна; $f, g$ сюръективны $\implies g \circ f$ сюръективна; $f, g$ биективны $\implies g \circ f$ биективна.\\
2. Обратные импликации: $g \circ f$ инъективна $\implies f$ инъективна; $g \circ f$ сюръективна $\implies g$ сюръективна.\\
3. Обращение: $(g \circ f)^{-1} = f^{-1} \circ g^{-1}$.
}
\ProofBox{%
Пусть $g \circ f$ инъективна. Если $f(a_1) = f(a_2)$, то $g(f(a_1)) = g(f(a_2)) \iff (g \circ f)(a_1) = (g \circ f)(a_2) \implies a_1 = a_2$, значит $f$ инъективна.
}

\ThmBox{Теорема 7.3 (Алгебра образов и прообразов)}{%
Для $f: A \to B$, подмножеств $A_i \subseteq A$ и $B_i \subseteq B$:
\begin{itemize}
\item $f(A_1 \cup A_2) = f(A_1) \cup f(A_2)$; но $f(A_1 \cap A_2) \subseteq f(A_1) \cap f(A_2)$ (равенство $\iff f$ инъективна!).
\item \textbf{Прообраз сохраняет абсолютно все теоретико-множественные операции}:
\[ f^{-1}(B_1 \cup B_2) = f^{-1}(B_1) \cup f^{-1}(B_2), \quad f^{-1}(B_1 \cap B_2) = f^{-1}(B_1) \cap f^{-1}(B_2), \quad f^{-1}(B_1 \setminus B_2) = f^{-1}(B_1) \setminus f^{-1}(B_2). \]
\item Монотонность: $X \subseteq Y \implies f(X) \subseteq f(Y)$ и $U \subseteq V \implies f^{-1}(U) \subseteq f^{-1}(V)$.
\end{itemize}
}
\ProofBox{%
$a \in f^{-1}(B_1 \cap B_2) \iff f(a) \in B_1 \cap B_2 \iff f(a) \in B_1 \land f(a) \in B_2 \iff a \in f^{-1}(B_1) \land a \in f^{-1}(B_2) \iff a \in f^{-1}(B_1) \cap f^{-1}(B_2)$. Равносильности строгие благодаря однозначности функции.
}

\ThmBox{Теорема 7.4 (Принцип Дирихле и пределы сжатия данных)}{%
Если $|A| > |B|$, ни одна функция $f: A \to B$ не может быть инъективной.
Следствие: никакой алгоритм сжатия без потерь не может сжимать все файлы (число строк длины $< n$ строго меньше $256^n$).
}

\ThmBox{Теорема 7.5 (Канторова функция спаривания)}{%
Функция $J: \mathbb{N}_0 \times \mathbb{N}_0 \to \mathbb{N}_0$, $J(x, y) = \frac{(x + y)(x + y + 1)}{2} + x$ является \textbf{биекцией}.
}

\DefBox{Определение 7.3 (Асимптотические классы Ландау)}{%
$f(n) = \mathcal{O}(g(n)) \iff \exists C, n_0 : \forall n \ge n_0 : f(n) \le C g(n)$; $f(n) = \Omega(g(n)) \iff \exists c, n_0 : f(n) \ge c g(n)$; $f(n) = \Theta(g(n)) \iff f \in \mathcal{O}(g) \cap \Omega(g)$; $f(n) = o(g(n)) \iff \lim \frac{f(n)}{g(n)} = 0$; $f(n) = \omega(g(n)) \iff \lim \frac{f(n)}{g(n)} = \infty$.
}

\section{Мощность множеств, счётность и континуум (Лекция 8)}

\DefBox{Определение 8.1 (Равномощность и бесконечность по Дедекинду)}{%
$A \sim B \iff$ существует биекция $f: A \leftrightarrow B$. Отношение $\sim$ --- эквивалентность, его классы --- кардинальные числа.\\
Множество $A$ \textbf{бесконечно по Дедекинду} $\iff \exists B \subsetneq A : B \sim A$ (равномощно собственному подмножеству).
}

\PropBox{Свойство 8.1 (Отель Гильберта и счётные множества)}{%
$\aleph_0 + 1 = \aleph_0$ ($n \mapsto n+1$); $\quad \aleph_0 + \aleph_0 = \aleph_0$ (чётные и нечётные); $\quad \aleph_0 \cdot \aleph_0 = \aleph_0$ (пары степенями простых $2^n 3^m$ или диагональной змейкой).\\
Множество счётно, если $A \sim \mathbb{N}$ ($|A| = \aleph_0$). Счётность --- наименьшая бесконечная мощность (всякое бесконечное подмножество $\mathbb{N}$ счётно).
}

\ThmBox{Теорема 8.1 (Счётность числовых множеств)}{%
Множества целых чисел $\mathbb{Z}$, рациональных чисел $\mathbb{Q}$ и алгебраических чисел $\mathbb{A}$ счётны.
}
\ProofBox{%
Для $\mathbb{A}$: число $\alpha$ --- корень $a_n x^n + \dots + a_0 = 0$ ($a_i \in \mathbb{Z}$). Высота $H = n + |a_n| + \dots + |a_0|$. При каждом фиксированном $H \ge 2$ число многочленов конечно, каждый имеет не более $n$ корней. Множество $\mathbb{A} = \bigcup_{H=2}^\infty K_H$ есть счётное объединение конечных множеств $\implies |\mathbb{A}| = \aleph_0$.
}

\ThmBox{Теорема 8.2 (Диагональный метод Кантора: несчётность $\mathbb{R}$)}{%
Отрезок $[0, 1]$ и множество всех вещественных чисел $\mathbb{R}$ \textbf{несчётны}: $|\mathbb{R}| = \mathfrak{c} = 2^{\aleph_0} > \aleph_0$.
}
\ProofBox{%
Пусть $[0, 1)$ счётен: $x_1, x_2, \dots$ Запишем десятичные знаки: $x_i = 0.d_{i1} d_{i2} d_{i3} \dots$ Построим $x^* = 0.c_1 c_2 \dots$, где $c_n = 4$, если $d_{nn} \ne 4$, и $c_n = 5$, если $d_{nn} = 4$. Число $x^* \in [0, 1)$, но $x^* \ne x_n$ для каждого $n$ (отличается $n$-я цифра). Числа $x^*$ нет в списке $\implies$ противоречие.
}

\ThmBox{Теорема 8.3 (Равномощность отрезка и квадрата / Совершенное множество Кантора)}{%
1. $[0, 1] \sim [0, 1]^2 \sim \mathbb{R}^n \sim \mathfrak{c}$ (доказывается переплетением цифровых блоков без девяток или кривой Пеано).\\
2. Канторово множество $\mathcal{C} = \bigcap C_n$ имеет меру Лебега $\mu(\mathcal{C}) = 0$, но его мощность равна континууму $|\mathcal{C}| = 2^{\aleph_0} = \mathfrak{c}$ (троичные дроби только из цифр 0 и 2).\\
3. Множество трансцендентных чисел $\mathbb{T} = \mathbb{R} \setminus \mathbb{A}$ имеет мощность континуума $|\mathbb{T}| = \mathfrak{c} - \aleph_0 = \mathfrak{c}$.
}

\ThmBox{Теорема 8.4 (Теорема Кантора о мощности булеана)}{%
Для любого множества $A$: $|A| < |\mathcal{P}(A)|$. Максимальной мощности не существует.
}
\ProofBox{%
Инъекция $x \mapsto \{x\}$ тривиальна. Покажем невозможность сюръекции $f: A \twoheadrightarrow \mathcal{P}(A)$.
Пусть $f$ сюръективна. Рассмотрим диагональное множество $B = \{x \in A \mid x \notin f(x)\} \subseteq A$.
Существует прообраз $b \in A$ такой, что $f(b) = B$.
Если $b \in B$, то по условию $b \notin f(b) = B$ (противоречие). Если $b \notin B$, то условие выполнено $\implies b \in f(b) = B$ (противоречие). Сюръекции нет $\implies |A| < |\mathcal{P}(A)|$.
}

\ThmBox{Теорема 8.5 (Кантора --- Бернштейна --- Шрёдера)}{%
Если $|A| \le |B|$ и $|B| \le |A|$, то $|A| = |B|$.
}
\ProofBox{%
Пусть даны инъекции $f: A \to B$ и $g: B \to A$. Проследим предков элементов при чередовании $f^{-1}$ и $g^{-1}$. Множество $A$ распадается на $A_A$ (цепи с началом в $A$, циклы и бесконечные цепи) и $A_B$ (цепи с началом в $B$). Биекция: $h(x) = f(x)$ при $x \in A_A$, и $h(x) = g^{-1}(x)$ при $x \in A_B$.
}

\PropBox{Свойство 8.2 (Шкалы Алефов и Бетов, Континуум-гипотеза)}{%
Алефы: $\aleph_0 < \aleph_1 < \aleph_2 < \dots$; Беты: $\beth_0 = \aleph_0, \beth_1 = 2^{\aleph_0} = \mathfrak{c}, \beth_{\alpha+1} = 2^{\beth_\alpha}$.
Континуум-гипотеза (CH): $\aleph_1 = \mathfrak{c}$. \textbf{Независима от ZFC} (Гёдель 1938, Коэн 1963).
}

\section{Вполне упорядоченные множества, ординалы и аксиоматика ZFC (Лекция 9)}

\DefBox{Определение 9.1 (Вполне упорядоченное множество / WOS)}{%
Линейный порядок $(X, \le)$ называется \textbf{вполне упорядоченным}, если любое непустое $S \subseteq X$ имеет наименьший элемент.
Критерий фундированности: в $(X, \le)$ нет бесконечной строго убывающей последовательности $a_0 \succ a_1 \succ a_2 \succ \dots$
}

\ThmBox{Теорема 9.1 (Сравнение вполне упорядоченных множеств)}{%
Для любых двух WOS $A$ и $B$ верно ровно одно из трёх: 1) $A \cong B$; 2) $A$ изоморфно начальному отрезку $B_{< b}$; 3) $B$ изоморфно начальному отрезку $A_{< a}$. Никакое WOS не изоморфно своему собственному начальному отрезку.
}

\DefBox{Определение 9.2 (Ординалы по фон Нейману)}{%
Множество $T$ называется \textbf{транзитивным}, если $\forall x \in T: x \subseteq T$.
Множество $\alpha$ называется \textbf{ординалом}, если $\alpha$ транзитивно и вполне упорядочено отношением $\in$.
Свойства: $\beta < \alpha \iff \beta \in \alpha$, а $\beta \le \alpha \iff \beta \subseteq \alpha$.
Числа фон Неймана: $0 = \varnothing, 1 = \{0\}, 2 = \{0, 1\}, n+1 = n \cup \{n\}$. Непосредственно следующий ординал: $\alpha^+ = \alpha \cup \{\alpha\}$.
}

\DefBox{Определение 9.3 (Классификация ординалов и трансфинитная индукция)}{%
Ординал бывает: 1) нулём $0$; 2) изолированным $\alpha = \beta + 1$; 3) предельным: $\lambda = \sup_{\beta < \lambda} \beta = \bigcup_{\beta < \lambda} \beta$.
Первый бесконечный ординал: $\omega = \{0, 1, 2, \dots\} = \mathbb{N}_0$.
\textbf{Трансфинитная индукция:} Доказывается: 1) $P(0)$; 2) $P(\alpha) \implies P(\alpha+1)$; 3) $(\forall \beta < \lambda : P(\beta)) \implies P(\lambda)$. Тогда $P(\alpha)$ истинно для всех ординалов.
}

\PropBox{Свойство 9.1 (Арифметика ординалов)}{%
\begin{itemize}
\item Сложение некоммутативно: $1 + \omega = \omega \ne \omega + 1$ (добавление точки в начало даёт $\omega$, в конец --- создаёт максимум).
\item Умножение некоммутативно: $2 \cdot \omega = \omega \ne \omega \cdot 2 = \omega + \omega$.
\item Ординальная степень: $2^\omega = \sup_{n < \omega} 2^n = \omega$ (счётно в ординалах, в отличие от кардинала $2^{\aleph_0} = \mathfrak{c}$).
\item Начальные ординалы: $\omega_\alpha \iff \aleph_\alpha$. Первый несчётный ординал $\omega_1$ (множество всех счётных ординалов).
\end{itemize}
}

\PropBox{Свойство 9.2 (Парадоксы теории множеств и 10 аксиом ZFC)}{%
Парадоксы: Рассела ($R = \{x \mid x \notin x\}$), Кантора (множество всех множеств $V$, $|V| < |\mathcal{P}(V)|$), Бурали-Форти (множество всех ординалов $\mathrm{On}$).
\textbf{Аксиоматика ZFC:}
1. \textbf{Экстенсиональность:} множества равны, если состоят из одних элементов.
2. \textbf{Пустое множество:} существует $\varnothing$.
3. \textbf{Пара:} существует $\{a, b\}$.
4. \textbf{Объединение:} существует $\bigcup A$.
5. \textbf{Схема выделения:} $\{x \in A \mid P(x)\}$ --- выделение только из существующего $A$ (блокирует парадокс Рассела).
6. \textbf{Булеан:} существует $\mathcal{P}(A)$.
7. \textbf{Бесконечность:} существует индуктивное множество (гарантирует $\omega$).
8. \textbf{Схема подстановки (замены):} образ множества при функции является множеством.
9. \textbf{Регулярность (фундирование):} каждое непустое множество содержит $\in$-минимальный элемент (запрещает $x \in x$ и бесконечные цепочки $\dots \in x_2 \in x_1$).
10. \textbf{Аксиома выбора (AC):} декартово произведение непустых множеств непусто.
Кумулятивная иерархия: $V_0 = \varnothing, V_{\alpha+1} = \mathcal{P}(V_\alpha), V_\lambda = \bigcup_{\beta < \lambda} V_\beta$.
Завершимость программ: сопоставление состоянию ординала, строго убывающего на каждом шаге (отсутствие бесконечного спуска $\implies$ завершение цикла).
}

\section{Язык и истина: синтаксис, семантика и компактность (Лекция 10)}

\DefBox{Определение 10.1 (Правильно построенная формула / WFF)}{%
Грамматика WFF: 1) Атомы $p \in V$ --- формулы; 2) Если $\varphi$ формула, то $(\neg \varphi)$ формула; 3) Если $\varphi, \psi$ формулы, то $(\varphi \circ \psi)$ формула ($\circ \in \{\wedge, \vee, \to, \leftrightarrow\}$). Дерево разбора: листья --- атомы, корень --- \textbf{главная связка}.
}

\ThmBox{Теорема 10.1 (Баланс скобок и однозначность разбора)}{%
1. $\operatorname{bal}(\varphi) = 0$, и для любого собственного непустого префикса $u$: $\operatorname{bal}(u) > 0$.\\
2. Каждая формула обладает \textbf{ровно одним} деревом синтаксического разбора.\\
3. Связки и листья: в любой формуле число бинарных связок $b$ и листьев $l$ связано соотношением $b = l - 1$.
}
\ProofBox{%
Индукцией по построению: для атома $b=0, l=1 \implies 0 = 1 - 1$. Отрицание не меняет $b$ и $l$. Для $(\psi_1 \circ \psi_2)$: $b = b_1 + b_2 + 1 = (l_1 - 1) + (l_2 - 1) + 1 = (l_1 + l_2) - 1 = l - 1$. Однозначность: первый префикс после первой скобки с нулевым внутренним балансом однозначно определяет левый операнд $\psi_1$.
}

\ThmBox{Теорема 10.2 (Композициональность и сведение следования к невыполнимости)}{%
1. \textbf{Композициональность (Тарский --- Фреге):} значение $\nu(\varphi)$ зависит исключительно от значений переменных, входящих в $\varphi$.\\
2. \textbf{Сведение следования к невыполнимости:}
\[ \Gamma \models \psi \iff \Gamma \cup \{\neg \psi\} \text{ невыполнимо}. \]
}
\ProofBox{%
($\implies$) Если $\Gamma \models \psi$, то на любой модели $\nu \models \Gamma$ истинно $\nu(\psi) = 1 \implies \nu(\neg\psi) = 0$. Значит, совместной модели у $\Gamma \cup \{\neg\psi\}$ нет. ($\impliedby$) Если совместной модели нет, то везде, где $\nu \models \Gamma$, значение $\nu(\neg\psi)$ не может быть 1, значит $\nu(\psi) = 1$, т.е. $\Gamma \models \psi$.
}

\ThmBox{Теорема 10.3 (Семантическая теорема о дедукции и принцип подстановки)}{%
1. \textbf{Семантическая дедукция:} $\Gamma \cup \{\varphi\} \models \psi \iff \Gamma \models (\varphi \to \psi)$.\\
2. \textbf{Принцип подстановки:} если $\varphi \equiv \psi$, то замена вхождения $\varphi$ на $\psi$ в формуле $\chi$ даёт эквивалентную формулу: $\chi[\varphi] \equiv \chi[\psi]$.
}

\ThmBox{Теорема 10.4 (Функциональная полнота штриха Шеффера / NAND)}{%
Связка $p \uparrow q = \neg(p \land q)$ (NAND) функционально полна:
\[ \neg p \equiv p \uparrow p, \quad p \land q \equiv (p \uparrow q) \uparrow (p \uparrow q), \quad p \lor q \equiv (p \uparrow p) \uparrow (q \uparrow q), \quad p \to q \equiv p \uparrow (q \uparrow q). \]
}

\ThmBox{Теорема 10.5 (Теорема компактности Мальцева --- Гёделя)}{%
Множество формул $\Gamma$ выполнимо тогда и только тогда, когда выполнимо любое его \textbf{конечное} подмножество $\Delta \subseteq \Gamma$.
}
\ProofBox{%
($\implies$) Очевидно. ($\impliedby$) Занумеруем переменные $p_1, p_2, \dots$ Построим бинарное дерево частичных означиваний длины $n$. Узел $\sigma \in \{0, 1\}^n$ назовём хорошим, если любое конечное $\Delta \subseteq \Gamma$ имеет модель, продолжающую $\sigma$. Корень хорош. По принципу Дирихле у хорошего узла хотя бы один потомок хорош (иначе объединение двух контрпримеров $\Delta_0 \cup \Delta_1$ конечно и не выполнимо, противоречие). По \textbf{лемме Кёнига} (в бесконечном конечно ветвящемся дереве есть бесконечная ветвь) получаем тотальную оценку $\nu$. Любая $\varphi \in \Gamma$ конечна и зависит от $p_1, \dots, p_N$. На шаге $N$ узел хорош, значит $\{\varphi\}$ выполнимо согласованно с $\nu \implies \nu(\varphi) = 1$. Модель построена.
}

\ThmBox{Теорема 10.6 (Конечные свидетели и теорема де Брёйна --- Эрдёша)}{%
1. \textbf{Конечные свидетели:} $\Gamma \models \psi \implies \exists \Gamma_0 \subseteq \Gamma, |\Gamma_0| < \infty : \Gamma_0 \models \psi$.\\
2. \textbf{Де Брёйн --- Эрдёш (1951):} бесконечный граф $k$-раскрашиваем $\iff$ любой его конечный подграф $k$-раскрашиваем (доказывается сведением к компактности логики высказываний).
}

\section{Натуральная дедукция и теория доказательств (Лекция 11)}

\DefBox{Определение 11.1 (Натуральная дедукция и нотация Фитча)}{%
Натуральная дедукция (Генцен, Яськовский, 1934) выводит теоремы через рассуждения от временных локальных гипотез с их последующим снятием (discharge). В нотации Фитча область действия гипотезы выделяется вертикальной чертой (subproof). Выход из subproof снимает гипотезу. Формулы внутри закрытого блока недоступны снаружи (лексическая область видимости).
}

\PropBox{Свойство 11.1 (Правила введения ($I$) и устранения ($E$) связок)}{%
\begin{itemize}
\item \textbf{Конъюнкция:} $\wedge I: \frac{A \quad B}{A \wedge B}; \qquad \wedge E_1: \frac{A \wedge B}{A}; \qquad \wedge E_2: \frac{A \wedge B}{B}$.
\item \textbf{Импликация:} $\to I: \frac{[A] \vdash B}{A \to B}$ (открытие гипотезы $A$, вывод $B$, закрытие subproof); $\qquad \to E: \frac{A \quad A \to B}{B}$ (Modus Ponens).
\item \textbf{Дизъюнкция:} $\vee I_1: \frac{A}{A \vee B}; \quad \vee I_2: \frac{B}{A \vee B}; \qquad \vee E: \frac{A \vee B \quad [A] \vdash C \quad [B] \vdash C}{C}$ (разбор случаев).
\item \textbf{Ложь и отрицание:} $\neg E: \frac{A \quad \neg A}{\bot}; \qquad \neg I: \frac{[A] \vdash \bot}{\neg A}; \qquad \bot E: \frac{\bot}{A}$ (Ex Falso Quodlibet / принцип взрыва).
\end{itemize}
}

\ThmBox{Теорема 11.1 (Классическая логика NK vs Интуиционистская NJ)}{%
Правила $\wedge, \lor, \to, \neg, \bot$ образуют интуиционистскую систему NJ. Классическая система NK получается добавлением одного из трёх равносильных правил:
1) \textbf{Косвенное доказательство (IP / Reductio ad Absurdum):} $\frac{[\neg A] \vdash \bot}{A}$;
2) \textbf{Снятие двойного отрицания (DNE):} $\frac{\neg\neg A}{A}$;
3) \textbf{Закон исключённого третьего (LEM):} $A \vee \neg A$.
\textbf{Принципиальное отличие $\neg I$ от IP:} правило $\neg I$ конструктивно (из предположения $A$ выводит противоречие и заключает $\neg A$); правило IP неконструктивно (из предположения $\neg A$ заключает бытие $A$ без предъявления конструкции).
}

\PropBox{Свойство 11.2 (Изоморфизм Карри --- Говарда)}{%
\[ \mathbf{Формулы \equiv Типы, \qquad Доказательства \equiv Программы, \qquad Нормализация \equiv \beta\text{-редукция}}. \]
$A \wedge B \equiv \text{тип пары } (A, B)$; $\quad A \to B \equiv \text{тип функции } A \to B$; $\quad A \vee B \equiv \text{тип суммы } \mathrm{Either}(A, B)$; $\quad \bot \equiv \text{пустой тип Void}$.
}

\ThmBox{Теорема 11.2 (Нормализация Правича и свойство подформульности)}{%
1. Всякое доказательство в натуральной дедукции приводится к нормальной форме без обходов (detours, введение связки с немедленным удалением) за конечное число редукций.\\
2. \textbf{Свойство подформульности:} в нормализованном выводе каждая встречающаяся формула является подформулой либо посылок $\Gamma$, либо заключения $\varphi$. Это делает поиск вывода в логике высказываний алгоритмически разрешимым.
}

\ThmBox{Теорема 11.3 (Корректность и Полнота натуральной дедукции)}{%
\[ \Gamma \vdash \varphi \iff \Gamma \models \varphi. \]
1. \textbf{Корректность (Soundness):} $\Gamma \vdash \varphi \implies \Gamma \models \varphi$ (индукцией по выводу: каждое правило сохраняет истинность).\\
2. \textbf{Лемма Линденбаума:} любое непротиворечивое множество формул содержится в максимальном непротиворечивом множестве $\Delta^*$.\\
3. \textbf{Лемма об истинности (Truth Lemma):} на канонической модели $\nu(p) = 1 \iff p \in \Delta^*$ для любой формулы $\psi$ выполнено $\nu(\psi) = 1 \iff \psi \in \Delta^*$.\\
4. \textbf{Полнота (Completeness):} $\Gamma \models \varphi \implies \Gamma \vdash \varphi$ (доказательство от противного через каноническую контрмодель).
}

\ThmBox{Теорема 11.4 (Теорема компактности как следствие полноты)}{%
Множество формул $\Gamma$ выполнимо $\iff$ каждое его конечное подмножество выполнимо.
}
\ProofBox{%
Пусть каждое конечное $\Gamma_0 \subseteq \Gamma$ выполнимо, но $\Gamma$ невыполнимо. Тогда $\Gamma \models \bot$. По полноте $\Gamma \vdash \bot$. Любой вывод состоит из конечного числа строк, в нём задействовано лишь конечное число посылок $\Gamma_0 \subset \Gamma$. Значит, $\Gamma_0 \vdash \bot$. По корректности $\Gamma_0 \models \bot$, то есть $\Gamma_0$ невыполнимо. Противоречие. Значит, $\Gamma$ выполнимо.
}

\section{Логика первого порядка / First-Order Logic (Лекция 11 репозитория)}

\DefBox{Определение 12.1 (Сигнатура, термы и формулы первого порядка)}{%
\textbf{Сигнатура} $\Sigma = (\mathcal{F}, \mathcal{P}, \mathrm{ar})$ задаёт функциональные и предикатные символы с их местностями.
\begin{itemize}
\item \textbf{Термы:} переменные $x, y, \dots$ и константы $c$ --- термы; если $t_1, \dots, t_k$ --- термы и $f \in \mathcal{F}$ ($k$-местная), то $f(t_1, \dots, t_k)$ --- терм. Термы обозначают объекты домена.
\item \textbf{Атомарные формулы:} $P(t_1, \dots, t_m)$ для $P \in \mathcal{P}$, и $t_1 = t_2$.
\item \textbf{Формулы:} строятся из атомарных формул с помощью связок $\neg, \land, \lor, \to$ и кванторов $\forall x, \exists x$.
\end{itemize}
}

\DefBox{Определение 12.2 (Интерпретация, модели и истинность)}{%
\textbf{Интерпретация (структура)} $\mathcal{M} = (D, I)$ состоит из непустого домена $D$ и функции интерпретации $I$, сопоставляющей каждой константе элемент $D$, функции --- операцию $D^k \to D$, предикату --- отношение $D^m \to \{0, 1\}$.
При означивании свободных переменных $\sigma: \mathrm{Var} \to D$ формула $\varphi$ принимает значение истины: $\mathcal{M}, \sigma \models \varphi$.
Формула \textbf{общезначима} ($\models \varphi$), если она истинна во всех интерпретациях при любых означиваниях.
}

\ThmBox{Теорема 12.1 (Законы распределения кванторов первого порядка)}{%
1. $\forall x (\varphi(x) \land \psi(x)) \equiv (\forall x \varphi(x)) \land (\forall x \psi(x))$.\\
2. $\exists x (\varphi(x) \lor \psi(x)) \equiv (\exists x \varphi(x)) \lor (\exists x \psi(x))$.\\
3. \textbf{Кванторные ловушки (неравносильности!):}\\
$\forall x (\varphi(x) \lor \psi(x)) \not\equiv (\forall x \varphi(x)) \lor (\forall x \psi(x))$ (контрпример: $x$ чётно или нечётно на $\mathbb{N}$);\\
$\exists x (\varphi(x) \land \psi(x)) \not\equiv (\exists x \varphi(x)) \land (\exists x \psi(x))$.
}

\DefBox{Определение 12.3 (Пренексная форма и сколемизация)}{%
\begin{itemize}
\item \textbf{Пренексная нормальная форма (ПНФ):} формула вида $Q_1 x_1 \dots Q_k x_k \, M$, где все кванторы $Q_i \in \{\forall, \exists\}$ вынесены вперёд, а матрица $M$ бескванторна. Всякая формула первого порядка алгоритмически приводится к логически эквивалентной ПНФ.
\item \textbf{Сколемизация:} замена квантора существования $\exists y$ в контексте $\forall x_1 \dots \forall x_m \exists y \, \psi$ на сколемовскую функцию $y = f(x_1, \dots, x_m)$ (или константу Сколема, если внешних $\forall$ нет). Результирующая формула \textbf{сколемовской нормальной формы (СНФ)} равновыполнима исходной.
\end{itemize}
}

\ThmBox{Теорема 12.2 (Теорема Чёрча о неразрешимости, 1936)}{%
Проблема общезначимости (и выполнимости) формул логики первого порядка \textbf{алгоритмически неразрешима}. Не существует единого алгоритма, определяющего по произвольной формуле первого порядка, является ли она общезначимой.
}

\ThmBox{Теорема 12.3 (Теорема Гёделя о полноте логики первого порядка, 1930)}{%
В исчислении предикатов первого порядка синтаксическая выводимость в точности совпадает с семантическим следованием:
\[ \Gamma \models \varphi \iff \Gamma \vdash \varphi. \]
}

\ThmBox{Теорема 12.4 (Теорема компактности первого порядка)}{%
Множество предложений первого порядка $\Gamma$ имеет модель тогда и только тогда, когда каждое его конечное подмножество $\Delta \subseteq \Gamma$ имеет модель.
}

\ThmBox{Теорема 12.5 (Теорема Лёвенгейма --- Сколема и некатегоричность)}{%
1. \textbf{Нисходящая часть:} если теория первого порядка в счётном языке имеет бесконечную модель, то она имеет счётную модель ($\aleph_0$).\\
2. \textbf{Восходящая часть:} если теория имеет бесконечную модель, она имеет модели сколь угодно большой мощности $\kappa$.\\
3. \textbf{Некатегоричность:} никакая теория первого порядка с бесконечной моделью не может зафиксировать мощность своей модели с точностью до изоморфизма (следствие: существование нестандартных моделей арифметики Пеано и нестандартного анализа).
}

\end{document}
